\scanpage{94}
\mspage{100}

On dit que $Z$ est \uline{complètemt} $\mathbb D$-\uline{asphérique}, si $Z\longrightarrow e$ est une $\mathbb D$-équivalence 2-complète.

Ainsi, l'axiome Der 3 \unsure{a''|a'}) peut \deleted{\unsure{se|s'être}} reformuler, en disant que si $Z$ est $\mathrm{Hot}$-asphérique, alors $Z$ est complètement $\mathbb D$-asphérique. On peut alors donner une variante affaiblie, qui suffira pour donner une notion raisonnable de $\mathbb D$-fibration :

% slanted
\margin{Et de $=$, Der 3 b', en disant que si $f:X\to Y$ \deleted{de} $\unsure{\in|\leftarrow}\mathrm{Diag}$ est une $\mathrm{Hot}$-\unsure{équiv.}, alors $f$ est une $\mathbb D$-équivalence 2-complète.}

\uline{Der 3} \unsure{b''|a''}) Toute $Z$ de $\mathrm{Diag}$ qui a soit un objet final, soit un objet initial, est complètement $\mathbb D$-asphérique.

\uline{Définition} On dit que $f:X\to Y$ dans $\mathrm{Diag}$ est une $\mathbb D$-fibration, si pour

\[
Y''\unsure{\relbar|\to}Y'\unsure{\relbar|\to}Y,
\]

avec $Y''\to Y'$ \deleted{\unsure{comme}} une $\mathbb D$-équivalence 1 complète (p. ex. $Y'$ et $Y''$ complètement $\mathbb D$-asphériques), $X''=X\times_Y Y''\unsure{\relbar|\to}X'=X\times_Y Y'$ est une $\mathbb D$-équivalence.

Cela équiv. à \deleted{\unsure{aux}}\inserted{chacune des} conditions b), b') de prop. 2. (page 79), i.e. que $f$ soit universellement $\mathbb D$-fibrant à gauche, ou $f$ univ. $\mathbb D$-fibrant à droite (grâce à Der 3 b'').

% archive sheet number: 93
